\chapter{Testing and Early Validation}
\label{ch:testing}
% ============================================================

\section{Testing Strategy}
\label{sec:testing-strategy}
Testing at this stage targets the data generation path rather than the detection
model, on the principle that model results obtained from an unverified dataset are
uninterpretable. Three levels are applied: component checks on extraction logic,
integration checks on the simulation-to-dataset path, and empirical validation that
injected attacks produce the theoretically expected signature.

\section{Test Levels and Techniques}
\label{sec:testing-levels}
\begin{itemize}
  \item \textbf{Component.} Assertions within the extraction pipeline on window
        counts, absence of missing values, and per-zone label distribution.
  \item \textbf{Integration.} End-to-end execution from configuration to dataset,
        verifying row counts and schema conformance.
  \item \textbf{Empirical validation.} Comparison of observed attack signatures
        against the expected theoretical effect, across multiple seeds.
  \item \textbf{Reproducibility.} Repeated execution with identical seeds, comparing
        generated structure and labels.
\end{itemize}

\section{Test Cases}
\label{sec:test-cases}

\begin{longtable}{p{0.07\textwidth} p{0.3\textwidth} p{0.3\textwidth} p{0.25\textwidth}}
\caption{Test cases executed at the mid-report stage}\\
\toprule
\textbf{ID} & \textbf{Objective} & \textbf{Expected} & \textbf{Result} \\
\midrule
\endfirsthead
\toprule
\textbf{ID} & \textbf{Objective} & \textbf{Expected} & \textbf{Result} \\
\midrule
\endhead
TC-01 & Topology instantiates with all declared nodes & All nodes present and linked &
Pass \\
TC-02 & gPTP synchronises all time-aware nodes & Node clocks track a common trajectory
under benign operation & Pass \\
TC-03 & Seeded reproducibility & Identical seed reproduces identical structure and labels &
Pass \\
TC-04 & Delay attack produces measurable deviation & Targeted branch diverges from
siblings during the attack interval & Pass \\
TC-05 & Drift attack produces zone-localisable deviation & Targeted zone diverges from
siblings & \textbf{Not realisable; see Section~\ref{sec:edge-cases}} \\
TC-06 & Extraction produces expected window count & Windows $=$ runs $\times$ zones
$\times$ windows-per-run $= 15 \times 4 \times 79$ & Pass (4{,}740) \\
TC-07 & Zero-event windows are retained, not dropped & Window with no clock event
yields a record with \texttt{n\_events}$=0$ & Pass \\
TC-08 & Attack labels attributed to victim zone only & In an attack run targeting zone
$z$, only zone $z$ carries the attack label & Pass (see Section~\ref{sec:test-results}) \\
TC-09 & Dataset conforms to the documented schema & Column names, order and types match
Table~\ref{tab:schema} & Pass (after schema amendment; see Section~\ref{sec:defects}) \\
\bottomrule
\end{longtable}

\section{Test Results}
\label{sec:test-results}

Fifteen seeded runs produced 225 recorded vectors, from which the extraction pipeline
generated 4{,}740 labelled feature windows.

\begin{longtable}{p{0.14\textwidth} r r r}
\caption{Label distribution by zone}\label{tab:labels}\\
\toprule
\textbf{Zone} & \textbf{Benign} & \textbf{Attack} & \textbf{Total} \\
\midrule
\endfirsthead
\toprule
\textbf{Zone} & \textbf{Benign} & \textbf{Attack} & \textbf{Total} \\
\midrule
\endhead
Zone 0 & 1{,}044 & 141 & 1{,}185 \\
Zone 1 & 1{,}044 & 141 & 1{,}185 \\
Zone 2 & 1{,}044 & 141 & 1{,}185 \\
Zone 3 & 1{,}044 & 141 & 1{,}185 \\
\midrule
\textbf{Total} & \textbf{4{,}176} & \textbf{564} & \textbf{4{,}740} \\
\bottomrule
\end{longtable}

\textbf{Derivation of the distribution.} With a window of 0.25\,s, a stride of
0.125\,s and a simulation limit of 10\,s, each run yields 79 windows per zone. The
observed 1{,}185 windows per zone therefore correspond to exactly 15 runs. Attacks
begin at $t=4.0$\,s, so 47 of the 79 windows per run fall within the attack interval.
The observed 141 attack-labelled windows per zone correspond to $3 \times 47$, that is,
each zone is the designated victim in exactly three runs. The run set is accordingly
3 benign runs and 12 delay-attack runs, covering four victim zones at three seeds each,
giving $12 \times 47 = 564$ attack windows in total.

\textbf{Victim rotation.} Rotating the victim zone across runs is a deliberate design
choice. Had a single zone been targeted throughout, zone identity would correlate
perfectly with the label and any classifier could reach high accuracy by learning zone
index alone. Because every zone appears as victim in three runs and as bystander in
nine, zone identity carries no label information and attribution must be learned from
the timing features.

\textbf{Class imbalance.} The benign-to-attack ratio is approximately $7.4{:}1$.
Accuracy is therefore excluded as a reported measure; evaluation uses per-class
precision, recall and F1, macro-F1, and detection rate at a fixed false-positive rate.

\section{Acceptance Testing / Validation}
\label{sec:acceptance-validation}

Acceptance checks defined in Section~\ref{sec:acceptance-criteria} have been exercised for FR-01 through FR-09.
Checks for FR-10 to FR-15 and NFR-08 require the federation and
monitoring-interference work scheduled for Iteration~04 and are reported as planned in
Section~\ref{sec:validation-plan}.

Success criterion SC-2, that protocol-native features outperform features computable
without knowledge of network structure, has been tested centrally ahead of schedule,
because it is the premise on which the federated work depends and a null result would
have required repositioning the project. The experiment is a feature ablation with
identical models and identical splits across three feature sets:

\begin{itemize}
  \item \textbf{within\_zone} (24 features): all \texttt{ctrl\_} and \texttt{zone\_}
        statistics. These describe one zone's clock behaviour in isolation and require
        no knowledge of the synchronisation tree. This is the control arm.
  \item \textbf{residual} (2 features): the cross-zone sibling residuals only.
  \item \textbf{combined} (26 features): both.
\end{itemize}

Splits are formed by holding out whole runs, stratified so that both benign and attack
runs appear on each side, and averaged over five resplits. Window-level splitting is
not used: the extraction stride is half the window length, so adjacent windows share
half their samples and a random split would place near-duplicates on both sides.

\begin{longtable}{p{0.22\textwidth} p{0.14\textwidth} r r r}
\caption{Feature ablation, macro-F1 averaged over five run-level splits}
\label{tab:ablation}\\
\toprule
\textbf{Feature set} & \textbf{Model} & \textbf{Macro-F1} & \textbf{Attack recall} &
\textbf{Latency (windows)} \\
\midrule
\endfirsthead
\toprule
\textbf{Feature set} & \textbf{Model} & \textbf{Macro-F1} & \textbf{Attack recall} &
\textbf{Latency (windows)} \\
\midrule
\endhead
Majority-class floor & n/a & 0.468 & n/a & n/a \\
\midrule
within\_zone  & logistic reg. & 0.500 & 0.419 & 11.6 \\
within\_zone  & MLP           & 0.481 & 0.013 & 25.0 \\
within\_zone  & random forest & 0.536 & 0.259 & 12.7 \\
\midrule
residual      & logistic reg. & 0.951 & 1.000 & 2.0 \\
residual      & MLP           & 0.968 & 0.993 & 2.0 \\
residual      & random forest & \textbf{0.984} & 0.994 & 2.5 \\
\midrule
combined      & logistic reg. & 0.975 & 0.997 & 2.0 \\
combined      & MLP           & 0.966 & 0.919 & 2.0 \\
combined      & random forest & 0.969 & 0.930 & 2.0 \\
\bottomrule
\end{longtable}

\textbf{Result 1: within-zone features do not detect this attack.} All three
within-zone arms fall between 0.481 and 0.536 macro-F1 against a majority-class floor
of 0.468. The MLP arm attains 0.013 attack recall, that is, it classifies almost every
window as benign. Features describing a single zone's clock in isolation carry
essentially no information about this attack, because the attack does not make the
victim's clock behave abnormally in absolute terms; it makes the victim disagree with
its siblings.

\textbf{Result 2: two cross-zone features are sufficient.} The residual arm reaches
0.951--0.984 macro-F1 using two columns. The gain over the control arm is $+0.475$,
$+0.485$ and $+0.433$ macro-F1 for logistic regression, MLP and random forest
respectively. SC-2 is therefore met, and met by a margin that does not depend on model
choice.

\textbf{Result 3: the larger feature set is not the better one.} For the MLP and random
forest, the residual arm alone outperforms the combined arm (0.984 against 0.969 for
random forest). The 24 within-zone features dilute rather than augment the signal. This
is relevant to NFR-05: a detector intended to run on a resource-constrained zone
controller benefits from a deliberately small feature set, and the evidence here is that
the small set is also the more accurate one.

\textbf{Result 4: both residual features are necessary.} That logistic regression, a
linear decision boundary, reaches 0.951 macro-F1 on two features implies that the
classes are linearly separable in the two-dimensional residual space. They are not
separable in the offset dimension alone, as Section~\ref{sec:edge-cases} shows. The drift-slope residual
is therefore load-bearing despite appearing weak when summarised univariately, and the
mechanism by which it contributes is recorded in Section~\ref{sec:edge-cases}.

\textbf{Detection latency.} With a confirmation requirement of three consecutive
anomalous windows, detection fires two windows after attack onset for every
residual-based arm. At a stride of 0.125\,s and a synchronisation interval of
0.125\,s, that is two synchronisation intervals, which is the minimum achievable under
a three-window confirmation rule. The within-zone arms require 11.6 to 25.0 windows,
the upper value indicating that detection frequently does not occur at all within the
run.

\textbf{Error structure.} For the combined random-forest arm the confusion matrix over
a held-out split is $\left(\begin{smallmatrix}1389 & 3\\ 0 & 188\end{smallmatrix}\right)$
(rows true, columns predicted): no attack window is missed, and three benign windows
are misclassified. Logistic regression shows the same structure with nine false
positives and no false negatives. The error is therefore one-sided towards false alarm
rather than missed detection, which is the preferable direction for a safety-relevant
monitor, though it makes the false-alarm rate the quantity requiring most attention at
deployment scale (Section~\ref{sec:validation-plan}).

\section{Edge and Failure Cases}
\label{sec:edge-cases}

\textbf{Negative result: oscillator drift manipulation is not realisable as a runtime
attack under the current simulation framework.} An attack was designed in which the
scenario manager would alter the local oscillator drift rate of a victim node partway
through a run, producing a zone-localised divergence from the remaining zones.

The attack could not be implemented in that form. The clock module's drift-rate
parameter is not declared mutable in the INET version in use, so the scenario manager
cannot modify it at simulation time. The implementation therefore applies the
manipulated drift rate from $t=0$, which is why the \texttt{DriftAttack} configuration
appears in the extraction script's \texttt{ALWAYS\_ON\_CONFIGS} set rather than being
gated on the attack start time.

Applied from $t=0$, the manipulation is not an attack in any meaningful sense: it is a
node whose oscillator runs at a different constant rate for the whole simulation, which, in the clock-offset dimension measured here, is indistinguishable from legitimate
component variation and which the gPTP servo is specifically designed to compensate. The observed network-wide behaviour is the protocol
operating correctly, not a detection failure.

Three consequences are carried forward. First, the finding is a statement about attack
realisability under this framework, not about the detectability of drift attacks in
general; the stronger claim is not supported by this experiment and is not made.
Second, drift is removed as a committed target for the cross-zone offset residual,
reducing the committed attack set to delay injection and synchronisation starvation; a
separate candidate signal for constant drift is described below. Third, if a runtime drift attack
is required in FYP-2, it must be realised by a different mechanism, for example by
subclassing the clock module to expose a mutable rate, or by post-processing the
recorded trace, and the cost of that work is weighed against the alternative of
deepening the evaluation of the two realisable attack classes.

\textbf{Update from the ring topology.} Constant drift applied from $t=0$ was later
re-run on the five-zone ring with the gPTP correction signals recorded
(Section~\ref{sec:ring-results}). The clock offset still hides it, but
\texttt{gmRateRatio} moves from about 0.99999 to about 0.99800 at 2000~ppm, so the signal
that is absent from the offset dimension is present in the rate-ratio dimension. This does
not change the realisability finding above, since a runtime drift attack remains
unrealisable, and it leaves the committed attack set unchanged pending extraction of
rate-ratio features.

\textbf{Negative result: no fixed threshold separates attack from benign behaviour.}
Two tuned thresholds on the absolute sibling offset residual were evaluated as
non-learned baselines. The first was set to the largest residual observed on benign
training windows, which by construction admits all benign training behaviour. The
second was set to the 99.5th percentile of that distribution, discarding the extreme
tail. Both attained zero recall on held-out data: every attack window fell below both
thresholds.

The cause is the gPTP convergence transient. Table~\ref{tab:populations} separates
benign windows by simulation time.

\begin{longtable}{p{0.34\textwidth} r r}
\caption{Absolute sibling offset residual by population}\label{tab:populations}\\
\toprule
\textbf{Population} & \textbf{Median (s)} & \textbf{Maximum (s)} \\
\midrule
\endfirsthead
\toprule
\textbf{Population} & \textbf{Median (s)} & \textbf{Maximum (s)} \\
\midrule
\endhead
Benign, $t < 1$\,s (convergence) & $1.00\times10^{-8}$ & $1.51\times10^{-5}$ \\
Benign, $t \geq 1$\,s (steady state) & $4.83\times10^{-9}$ & $1.81\times10^{-6}$ \\
Delay attack & $2.19\times10^{-6}$ & $5.66\times10^{-6}$ \\
\bottomrule
\end{longtable}

Before the clock servos lock, zones disagree with one another by more than any injected
attack subsequently moves them: the convergence maximum exceeds the largest attack
residual by a factor of 2.7. Those early windows are legitimately benign, so any
threshold admitting all benign behaviour must sit above the startup excursion and is
consequently blind to the attack in steady state.

The percentile variant fails for a quantitative reason worth stating. Windows beginning
before $t = 1$\,s constitute approximately 11\% of all benign windows, so any percentile
above roughly the 89th remains inside the convergence population. The failure is
therefore not an artefact of a poorly chosen cut-off: no percentile-based threshold
succeeds while the convergence transient is part of the benign class.

This is a stronger justification for the learned detector than the ablation alone
provides. The ablation establishes which features carry the signal; this result
establishes that the signal is not separable by a single decision boundary on the
offset residual, because the benign class is not unimodal; it is a mixture of a
steady-state population and a convergence population, and the attack lies between them
in magnitude.

\textbf{Why two residual features are required.} The second residual feature,
computed on the drift slope rather than the mean offset, is a weak discriminator when
measured alone: the ratio of its median under attack to its median under benign
operation is approximately 2.0, against 438 for the offset residual. It is nevertheless
not redundant. The residual feature set contains only these two columns, and logistic
regression, a linear decision boundary, attains 0.951 macro-F1 on them
(Table~\ref{tab:ablation}). A linear boundary can only achieve this if the three
populations are linearly separable in the two-dimensional residual space, which they
are not in the offset dimension alone.

The mechanism is that the two populations the offset residual confuses are distinguished
by the slope. During convergence the servos are actively correcting, so successive
offsets change rapidly and the slope residual is large. The delay attack introduces a
steady asymmetry, so the victim's offset departs from its siblings along a near-constant
gradient and the slope residual remains small. The slope residual is therefore not a
second detector of the attack; it is the feature that excludes the benign convergence
population from the region of offset-residual space that the attack occupies. This is
the reason a univariate summary understates its contribution, and it is recorded here
because the measurement that reveals it, a linear model succeeding on two features,
is easily mistaken for an unimportant result.

The practical consequence for FYP-2 is that the attack-intensity sweep (Section~\ref{sec:validation-plan})
must report the benign envelope, separated into convergence and steady-state
populations, alongside detection rate rather than reporting detection rate alone.

\textbf{Pending: grandmaster-level perturbation.} A separate structural limitation
follows from the planned sibling-residual feature rather than from any implemented
result. A residual between siblings cancels common-mode perturbation by construction,
so it cannot detect an effect that moves the entire synchronisation tree together,
including compromise at the time source itself. This bounds the detection claim to
link-local timing integrity. Detecting time-source compromise would require an
independent reference such as local oscillator holdover comparison or a second
synchronisation domain, and is out of scope.

\textbf{Pending edge cases.} Attack intensities approaching the benign drift envelope,
unseen seeds, and coordinator or model unavailability are scheduled for Iteration~04
(Section~\ref{sec:validation-plan}).

\section{Defects and Corrective Actions}
\label{sec:defects}

\textbf{DEF-04. Schema drift between the documented column contract and the emitted
dataset.} Test TC-09 compared the columns of \texttt{features.csv} against the
contract in Table~\ref{tab:schema} and found three columns present in the data but
absent from the contract: \texttt{target\_zone}, and the two
\texttt{sibling\_n\_\{f\}} counts emitted alongside each residual.
\emph{Consequence:} a reader implementing against the documented contract would not
have anticipated them, and \texttt{target\_zone} in particular is equivalent to the
label, so an implementer unaware of it could leak ground truth into a model.
\emph{Recovery:} the contract in Table~\ref{tab:schema} was corrected to document all
three, and the leakage risk is additionally guarded by an automated test asserting
that no identifier column reaches any feature set (Appendix~\ref{app:test-cases}). Status:
\textbf{closed}.

\begin{longtable}{p{0.08\textwidth} p{0.3\textwidth} p{0.26\textwidth} p{0.28\textwidth}}
\caption{Defect register}\\
\toprule
\textbf{ID} & \textbf{Description} & \textbf{Impact} & \textbf{Corrective action} \\
\midrule
\endfirsthead
\toprule
\textbf{ID} & \textbf{Description} & \textbf{Impact} & \textbf{Corrective action} \\
\midrule
\endhead
DEF-01 & The feature set contained no cross-zone comparison; all features were
computed within a single zone. &
The features were not structurally exclusive to a deterministic network, so SC-2 could
not be tested. &
\texttt{build()} restructured to iterate window-then-zone, and two sibling residual
features emitted. Verified by the ablation in Section~\ref{sec:acceptance-validation}: within-zone features reach
0.48--0.54 macro-F1, residual features 0.95--0.98. Status: \textbf{closed}. \\
\addlinespace
DEF-02 & All four zones of the Iteration~1 topology carry identical traffic profiles
and produce identical window counts. &
Federated, centralized and local-only training would converge to equivalent models,
making the federated comparison (SC-4) null by construction. &
The five-zone reference topology (Section~\ref{sec:ref-topology}) introduces structural asymmetry: three
to seven ECUs per zone, distinct traffic classes by zone function, and gPTP delay that
grows with hop distance from the grandmaster. Status: \textbf{partially addressed}.
Designed but not measured; closes when per-zone divergence is quantified and
differentiated traffic profiles are running. \\
\addlinespace
DEF-03 & The attack start time is hard-coded in the extraction script as a module
constant and must be kept consistent with the simulation configuration by hand. &
A change to the attack schedule in the simulation that is not mirrored in the extraction
script would mislabel the dataset silently. &
Read the attack start time from run metadata rather than from a constant, and add an
assertion that fails extraction when the two disagree. Status: open, low effort. \\
\bottomrule
\end{longtable}

\section{Current Validation Status}

The data generation path is validated for one attack class and found defective in one
respect. The detection and federation paths are unvalidated because they are
unimplemented. No detection performance figures are claimed at this stage, and none
should be inferred from the dataset statistics reported above. The attack simulations on
the five-zone ring (Section~\ref{sec:ring-results}) are single-run observations of candidate
signals; they have not been extracted into the dataset, and no detection performance is
claimed from them.

% ============================================================